\documentclass[12pt,a4paper]{article}
\usepackage{ugmath}
\usepackage{float}
\usepackage{placeins}
\usepackage{array}
\usepackage[labelfont=bf,labelsep=space]{caption}
\newcommand{\paren}[1]{\left( #1 \right)}
\newcommand{\alumno}{Ricardo León Martínez}
\newcommand{\materia}{Algebra Lineal I}
\newcommand{\profesor}{Claudia Reynoso Alcántara}
\newcommand{\tarea}{Tarea 10}
\newcommand{\fecha}{27/4/2026}


\begin{document}

    \begin{center}
        {\large\textbf{UNIVERSIDAD DE GUANAJUATO}}\\[0.3cm]
        {\normalsize\textbf{DIVISIÓN DE CIENCIAS NATURALES Y EXACTAS}}\\
        {\normalsize\textbf{CAMPUS GUANAJUATO}}\\[1cm]

        {\Large\textbf{\tarea\ (\materia)}}\\[1cm]
    \end{center}

    \noindent
    \textbf{Nombre:} \alumno 
    \hfill 
    \textbf{Fecha:} \fecha 
    \hfill 
    \textbf{Calificación:} \rule{3cm}{0.4pt}

    \vspace{0.3cm}

    \begin{mdframed}[style=mdpinkbox,frametitle={Ejercicio 1}]
        (1 punto) Encuentra una base de $\Im(T)$ y de $\ker(T)$, donde $T:\R^{3}\to\R^{3}$ está
        dada por
        \begin{equation*}
            \begin{pmatrix}
                -1 & 2 & 1\\
                0 & 1 & 1\\
                1 & 3 & 4
            \end{pmatrix}
        \end{equation*}
    \end{mdframed}

    \begin{solucion}
        La imagen de $T$ es el subespacio generado por las columnas de $A$:
        \begin{equation*}
            \Im(T)=\text{span}\left\{
                \begin{pmatrix}
                    -1\\
                    0\\
                    1
                \end{pmatrix},
                \begin{pmatrix}
                    2\\
                    1\\
                    3
                \end{pmatrix},
                \begin{pmatrix}
                    1\\
                    1\\
                    4
                \end{pmatrix}
            \right\}
        \end{equation*}
        Veamos su independencia lineal reduciendo la matriz
        \begin{equation*}
            \begin{pmatrix}
                -1 & 2 & 1\\
                0 & 1 & 1\\
                1 & 3 & 4
            \end{pmatrix}
            \sim
            \begin{pmatrix}
                -1 & 2 & 1\\
                0 & 1 & 1\\
                0 & 5 & 5
            \end{pmatrix}
            \sim
            \begin{pmatrix}
                -1 & 2 & 1\\
                0 & 1 & 1\\
                0 & 0 & 0 
            \end{pmatrix}
        \end{equation*}
        ya que el rango de una matriz esta dado por el numero de pivotes se sigue que $\rg(T)=2$.
        Por tanto, dos columnas son linealmente independientes. Tomando las dos primeras columnas
        \begin{equation*}
            \BB_{\Im(T)}=\left\{
                \begin{pmatrix}
                    -1\\
                    0\\
                    1
                \end{pmatrix},
                \begin{pmatrix}
                    2\\
                    1\\
                    3
                \end{pmatrix}
            \right\}
        \end{equation*}
        Ahora, buscamos los vectores $x=\{x_{1},x_{2},x_{3}\}$ tales que $Ax=0$. El sistema es
        \begin{equation*}
            \begin{cases}
                -x_{1}+2x_{2}+x_{3}=0\\
                x_{2}+x_{3}=0\\
                x_{1}+3x_{2}+4x_{3}=0
            \end{cases}
        \end{equation*}
        De la segunda ecuación
        \begin{equation*}
            x_{2}=-x_{3}.
        \end{equation*}
        Sustituyendo en la primera
        \begin{equation*}
            -x_{1}+2(-x_{3})+x_{3}=x_{1}-2x_{3}=0\implies x_{1}=-x_{3}.
        \end{equation*}
        Por lo tanto
        \begin{equation*}
            x=x_{3}(-1,-1,1).
        \end{equation*}
        Luego
        \begin{equation*}
            \ker(T)=\text{span}\{(-1,-1,1)\}.
        \end{equation*}
        Así, una base es
        \begin{equation*}
            \BB_{\ker(T)}=\left\{
                \begin{pmatrix}
                    -1\\
                    -1\\
                    1
                \end{pmatrix}
            \right\}
        \end{equation*}
    \end{solucion}

    \begin{mdframed}[style=mdpinkbox,frametitle={Ejercicio 2}]
        (1 punto) Sea $V_{2}$ el espacio vectorial sobre $F$, de polinomios de grado $\leq2$ y
        considera la transformación lineal:
        \begin{align*}
            T:V_{2}&\to F^{3}\\
            a_{0}+a_{1}x+a_{2}x^{2}&\mapsto(a_{0},a_{1},a_{2}).
        \end{align*}
        Sean $\BB_{1}=\{1,x,x^{2}\}$ y $\BB_{2}=\{e_{1},e_{2},e_{3}\}$ bases ordenadas de $V_{2}$
        y $\R^{3}$. Encuentra $[T]_{\BB_{1},\BB_{2}}$
    \end{mdframed}

    \begin{solucion}
        Por definición,
        \begin{equation*}
            [T]_{\BB_{1},\BB_{2}}=([T(1)]_{\BB_{2}}\quad[T(x)]_{\BB_{2}}\quad[T(x^{2})]_{\BB_{2}}).
        \end{equation*}
        Calculamos
        \begin{equation*}
            T(1)=(1,0,0),\quad T(x)=(0,1,0),\quad T(x^{2})=(0,0,1).
        \end{equation*}
        Dado que $\BB_{2}$ es la base canónica, estas ya son sus coordendas. Así
        \begin{equation*}
            [T]_{\BB_{1},\BB_{2}}=
            \begin{pmatrix}
                1 & 0 & 0\\
                0 & 1 & 0\\
                0 & 0 & 1
            \end{pmatrix}.
        \end{equation*}
    \end{solucion}

    \begin{mdframed}[style=mdpinkbox,frametitle={Ejercicio 3}]
        (2 puntos) Sea $V=\{f:\R\to\R:f(x)=\sum_{i=0}^{3}a_{i}x^{i},a_{i}\in\R\}$. Considera la
        transformación lineal:
        \begin{align*}
            \TT:V&\to V,\\
            f(x):&\mapsto xf^{\prime}(x).
        \end{align*}
        Encuentra $[\TT]_{\BB_{1}}$, $[\TT]_{\BB_{2}}$, $[\TT]_{\BB_{1},\BB_{2}}$ y
        $[\TT^{2}]_{\BB_{1},\BB_{2}}$, donde $\BB_{1}=\{1,x,x^{2},x^{3}\}$ y $\BB_{2}=\{1,x+1,(x+1)^{2},(x+1)^{3}\}$.
    \end{mdframed}

    \begin{solucion}
        1. Matriz $[\TT]_{\BB_{1}}$ con $\BB_{1}=\{1,x,x^{2},x^{3}\}$\\
        Calculamos la acción sobre la base
        \begin{equation*}
            \TT(1)=0,\quad\TT(x)=x,\quad\TT(x^{2})=2x^{2},\quad\TT(x^{3})=3x^{3}.
        \end{equation*}
        Coordenadas en $\BB_{1}$
        \begin{equation*}
            [\TT]_{\BB_{1}}=
            \begin{pmatrix}
                0 & 0 & 0 & 0\\
                0 & 1 & 0 & 0\\
                0 & 0 & 2 & 0\\
                0 & 0 & 0 & 3
            \end{pmatrix}.
        \end{equation*}
        2. Matriz $[\TT]_{\BB_{2}}$, con $\BB_{2}=\{1,x+1,(x+1)^{2},(x+1)^{3}\}$\\
        Sea $y=x+1$, entonces $x=y-1$. Calculamos
        \begin{align*}
            \TT(1)&=0,\\
            \TT(x+1)&=x=(y-1)=(x+1)-1,\\
            \TT((x+1)^{2})&=x\cdot2(x+1)=2x(x+1)=2(y-1)y=2y^{2}-2y,\\
            \TT((x+1)^{3})&=x\cdot3(x+1)^{2}=3(y-1)y^{2}=3y^{3}-3y^{2}.
        \end{align*}
        Expresando en $\BB_{2}=\{1,y,y^{2},y^{3}\}$
        \begin{equation*}
            [\TT]_{\BB_{2}}=
            \begin{pmatrix}
                0 & -1 & 0 & 0\\
                0 & 1 & -2 & 0\\
                0 & 0 & 2 & -3\\
                0 & 0 & 0 & 3
            \end{pmatrix}
        \end{equation*}
        3. Matriz $[\TT]_{\BB_{1},\BB_{2}}$\\
        Matriz de cambio de base de $\BB_{1}$ a $\BB_{2}$
        \begin{equation*}
            1=1,\quad x+1=x+1\quad (x+1)^{2}=x^{2}+2x+1,\quad (x+1)^{3}=x^{3}+3x^{2}+3x+1.
        \end{equation*}
        Por tanto
        \begin{equation*}
            P=[I]_{\BB_{2}\to\BB_{1}}=
            \begin{pmatrix}
                1 & 1 & 1 & 1\\
                0 & 1 & 2 & 3\\
                0 & 0 & 1 & 3\\
                0 & 0 & 0 & 1
            \end{pmatrix}.
        \end{equation*}
        Por definición
        \begin{equation*}
            [\TT]_{\BB_{1},\BB_{2}}=P^{-1}[\TT]_{\BB_{1}}P.
        \end{equation*}
        Así,
        \begin{equation*}
            [\TT]_{\BB_{1},\BB_{2}}=
            \begin{pmatrix}
                1 & -1 & 1 & -1\\
                0 & 1 & -2 & 3\\
                0 & 0 & 1 & -3\\
                0 & 0 & 0 & 1
            \end{pmatrix}
            \begin{pmatrix}
                0 & 0 & 0 & 0\\
                0 & 1 & 0 & 0\\
                0 & 0 & 2 & 0\\
                0 & 0 & 0 & 3
            \end{pmatrix}
            \begin{pmatrix}
                1 & 1 & 1 & 1\\
                0 & 1 & 2 & 3\\
                0 & 0 & 1 & 3\\
                0 & 0 & 0 & 1
            \end{pmatrix}
        \end{equation*}
        Concluyendo que
        \begin{equation*}
            [\TT]_{\BB_{1},\BB_{2}}=
            \begin{pmatrix}
                0 & -1 & 0 & 0\\
                0 & 1 & -2 & 0\\
                0 & 0 & 2 & -3\\
                0 & 0 & 0 & 3
            \end{pmatrix}
        \end{equation*}
        4. $[\TT^{2}]_{\BB_{1},\BB_{2}}$\\
        Notemos que
        \begin{equation*}
            \TT(x^{n})=nx^{x}\implies\TT^{2}(x^{n})=n^{2}x^{n}.
        \end{equation*}
        Luego
        \begin{equation*}
            [\TT^{2}]_{\BB_{1}}=
            \begin{pmatrix}
                0 & 0 & 0 & 0\\
                0 & 1 & 0 & 0\\
                0 & 0 & 4 & 0\\
                0 & 0 & 0 & 9
            \end{pmatrix}.
        \end{equation*}
        y por cambio de base
        \begin{equation*}
            [\TT^{2}]_{\BB_{1},\BB_{2}}=P^{-1}[\TT^{2}]_{\BB_{1}}P.
        \end{equation*}
        Así,
        \begin{equation*}
            [\TT^{2}]_{\BB_{1},\BB_{2}}=
            \begin{pmatrix}
                1 & 0 & 3 & 1\\
                0 & 1 & -6 & 6\\
                0 & 0 & 4 & -15\\
                0 & 0 & 0 & 9
            \end{pmatrix}
        \end{equation*}
    \end{solucion}

    \begin{mdframed}[style=mdpinkbox,frametitle={Ejercicio 4}]
        (2 puntos) Determina si, sobre $\R$, las siguientes matrices son semejantes:
        \begin{equation*}
            A=
            \begin{pmatrix}
                4 & 1 & 0\\
                0 & 4 & 0\\
                0 & 0 & 2
            \end{pmatrix},
            \quad
            B=
            \begin{pmatrix}
                4 & 0 & 0\\
                0 & 4 & 1\\
                0 & 0 & 2
            \end{pmatrix}.
        \end{equation*}
    \end{mdframed}

    \begin{mdframed}[style=mdpinkbox,frametitle={Ejercicio 5}]
        (2 puntos) Sean $V$ y $W$, espacios vectoriales sobre un campo $F$, definimos el conjunto:
        \begin{equation*}
            \LL(V,W)=\{T:V\to W: T\text{ es lineal}\},
        \end{equation*}
        demuestra que, con las siguientes operaciones $\LL(V,W)$ tiene estructura de espacio
        vectorial sobre $F$:
        \begin{align*}
            +:\LL(V,W)\times\LL(V,W)&\to\LL(V,W)\\
            (T_{1},T_{2})&\mapsto T_{1}+T_{2}
        \end{align*}
        \begin{align*}
            \cdot:F\times\LL(V,W)&\to\LL(V,W)\\
            (a,T)&\mapsto aT
        \end{align*}
    \end{mdframed}

    \begin{proof}
        Sean $T_{1},T_{2}\in\LL(V,W)$. Entonces para $u,v\in V$ y $\lambda\in F$
        \begin{align*}
            (T_{1}+T_{2})(u+v)=T_{1}(u+v)+T_{2}(u+v)=T_{1}(u)+T_{1}(v)+T_{2}(u)+T_{2}(v)\\
            =(T_{1}(u)+T_{2}(u))+(T_{1}(v)+T_{2}(v))=(T_{1}+T_{2})(u)+(T_{1}+T_{2})(v),\\
            (T_{1}+T_{2})(\lambda u)=T_{1}(\lambda u)+T_{2}(\lambda u)=\lambda T_{1}(u)+\lambda T_{2}(u)=\lambda(T_{1}+T_{2})(u).
        \end{align*}
        Luego $T_{1}+T_{2}\in\LL(V,W)$.\\
        Sea $a\in F$, $T\in\LL(V,W)$
        \begin{align*}
            (aT)(u+v)=aT(u+v)=a(T(u)+T(v))=aT(u)+aT(v),\\
            (aT)(\lambda u)=aT(\lambda u)=a(\lambda T(u))=\lambda(aT(u)).
        \end{align*}
        Luego $aT\in\LL(V,W)$.\\
        Sean $T_{1},T_{2},T_{3}\in\LL(V,W)$, $a,b\in F$, $v\in V$.\\
        (i) Asociatividad de la suma
        \begin{equation*}
            ((T_{1}+T_{2})+T_{3})(v)=(T_{1}+T_{2})(v)+T_{3}(v)=T_{1}(v)+T_{2}(v)+T_{3}(v)=(T_{1}+(T_{2}+T_{3}))(v).
        \end{equation*}
        (ii) Conmutatividad
        \begin{equation*}
            (T_{1}+T_{2})(v)=T_{1}(v)+T_{2}(v)=T_{2}(v)+T_{1}(v)=(T_{2}+T_{1})(v).
        \end{equation*}
        (iii) Elemento neutro\\
        Definimos $0:V\to W$, $0(v)+0_{W}$. Entonces
        \begin{equation*}
            (T+0)(v)=T(v)+0_{W}=T(v).
        \end{equation*}
        (iv) Inverso aditivo\\
        Para $T$, definimos $(-T)(v):=-T(v)$. Entonces
        \begin{equation*}
            (T+(-T))(v)=T(v)-T(v)=0_{W}.
        \end{equation*}
        (v) Asociativida con escalares
        \begin{equation*}
            ((ab)T)(v)=(ab)T(v)=a(bT(v))=(a(bT))(v).
        \end{equation*}
        (vi) Distributividad
        \begin{align*}
            ((a+b)T)(v)=(a+b)T(v)=aT(v)+bT(v)=(aT+bT)(v),\\
            (a(T_{1}+T_{2}))(v)=a(T_{1}(v)+T_{2}(v))=aT_{1}(v)+aT_{2}(v).
        \end{align*}
        (vii) Neutro multiplicativo
        \begin{equation*}
            (1T)(v)=1\cdot T(v)=T(v).
        \end{equation*}
        Así, $\LL(V,W)$ es un espacio vectorial sobre $F$.
    \end{proof}
    
    \begin{mdframed}[style=mdpinkbox,frametitle={Ejercicio 6}]
        (1 punto) Demuestra que todo elemento en $(F^{n})^{*}$ es de la forma
        \begin{align*}
            F^{n}&\to F,\\
            (x_{1},\dots,x_{n})&\mapsto\sum_{i=1}^{n}a_{i}x_{i}
        \end{align*}
        donde $a_{1},\dots,a_{n}\in F$.
    \end{mdframed}

    \begin{proof}
        Sea $\{e_{1},\dots,e_{n}\}$ la base canónica de $F^{n}$, donde
        \begin{equation*}
            e_{i}=(0,\dots,0,1,0,\dots,0).
        \end{equation*}
        Todo vector $x\in F^{n}$ se escribe de manera unica como
        \begin{equation*}
            x=\sum_{i=1}^{n}x_{i}e_{i}.
        \end{equation*}
        Por linealidad de $T$
        \begin{equation*}
            T(x)=T\left(\sum_{i=1}^{n}x_{i}e_{i}\right)=\sum_{i=1}^{n}x_{i}T(e_{i}).
        \end{equation*}
        Definimos
        \begin{equation*}
            a_{i}:=T(e_{i})\in F.
        \end{equation*}
        Entonces
        \begin{equation*}
            T(x_{1},\dots,x_{n})=\sum_{i=1}^{n}a_{i}x_{i}.
        \end{equation*}
    \end{proof}

    \begin{mdframed}[style=mdpinkbox,frametitle={Ejercicio 7}]
        (1 punto) Considera la base ordenada $\BB=\{(1,0,-1),(1,1,1),(2,2,0)\}$ de $\R^{3}$.
        Encuentra la base dual.
    \end{mdframed}

\end{document}